\documentclass[11pt]{article}
\usepackage[margin=1in]{geometry}
\usepackage[T1]{fontenc}
\usepackage[utf8]{inputenc}
\usepackage{lmodern,microtype,amsmath,amssymb,amsthm,mathtools}
\usepackage[colorlinks=true,linkcolor=blue,citecolor=blue,urlcolor=blue]{hyperref}
\hypersetup{pdftitle={Symmetric parity predictors require large amplitudes},pdfauthor={Dmitry Vaintrob; Kaarel Hanni; OpenAI Codex}}
\setlength{\emergencystretch}{2em}
\newtheorem{theorem}{Theorem}
\newtheorem{lemma}[theorem]{Lemma}
\newtheorem{corollary}[theorem]{Corollary}
\newtheorem{proposition}[theorem]{Proposition}
\theoremstyle{remark}\newtheorem{remark}[theorem]{Remark}
\newcommand{\E}{\mathbb E}\newcommand{\Prb}{\mathbb P}
\newcommand{\R}{\mathbb R}\newcommand{\N}{\mathbb N}
\DeclareMathOperator{\sign}{sign}
\title{Symmetric parity predictors require large amplitudes\\
\large An exponential lower bound and a consequence for margin repairs}
\author{Dmitry Vaintrob \and Kaarel H\"anni \and OpenAI Codex}
\date{October 16, 2026}
\begin{document}
\maketitle
\begin{abstract}
A symmetric polynomial on the Boolean cube that predicts parity with
error at most $1/3$, has absolute value at least one everywhere, and has
degree at most a sufficiently small fixed fraction of the dimension
must have $L^1$ norm at least $\exp(c\sqrt{k})$. The allowed degree
fraction can be any fixed $\alpha<1-\sqrt3/2$.
The proof uses a binomial generating function, an elementary
Schr\"odinger barrier estimate, and an averaged form of Jensen's formula. It includes arbitrary choices of
sign pattern and polynomial amplitudes. We also prove that every such
symmetric head has diverging average cost in an explicitly defined
cusp model of margin repairs, even when repairs may change unobserved
signs at intermediate radii. This cost conclusion extends to arbitrary
amplitudes whenever the head's sign pattern is symmetric in a sufficiently
large block. Sign symmetry and eventual agreement with the head remain
explicit hypotheses; no reduction of a general Bayesian posterior to
this structured problem is asserted.
\end{abstract}

\section{Statement and scope}
Write $\chi_k(x)=\prod_{i=1}^k x_i$ for parity on $\{-1,1\}^k$, equipped
with the uniform measure. A symmetric real polynomial $P$ depends only
on the Hamming weight $H=|\{i:x_i=-1\}|$. We use $P(h)$ for its value
on that layer and
\[
 w_k(h)=2^{-k}\binom{k}{h},\qquad
 \|P\|_p=\left(\sum_{h=0}^k w_k(h)|P(h)|^p\right)^{1/p}.
\]
Polynomial degree means multilinear Fourier degree on the cube. For a
symmetric function this is also the degree in its Krawtchouk expansion.
Its parity advantage is
\[
 c(P)=\E[\chi_k\sign P]=1-2\Prb(\sign P\ne\chi_k).
\]
All logarithms below are natural.

\begin{theorem}[Exponential amplitude lower bound]\label{thm:main}
Fix $0<\rho<1$ and a degree fraction $\alpha$ satisfying
\begin{equation}\label{eq:criterion}
 0<\alpha<1-\cos(\pi\rho/2).
\end{equation}
There are constants $c>0$ and $k_0$, depending only on $\alpha,\rho$,
with the following property. If $k\ge k_0$ and $P$ is symmetric with
\[
 \deg P\le\alpha k,\qquad |P(x)|\ge1\text{ for every }x,
 \qquad c(P)\ge\rho,
\]
then $\|P\|_1\ge\exp(c\sqrt{k})$.
\end{theorem}

For error at most $1/3$, take $\rho=1/3$. The allowed degree fraction
is any fixed $\alpha<1-\sqrt3/2=0.1339745962\ldots$.
The constants and the threshold are not claimed optimal.

The theorem concerns every symmetric polynomial of the indicated
fractional degree, not merely a particular interpolation polynomial or
a prescribed sign pattern. The norm conclusion is stronger than any
polynomial lower bound in $k$. It does not hold for unrestricted
nonsymmetric heads: substituting growing disjoint block parities into
a fixed accurate polynomial preserves its norm and relative degree.
Such an outer polynomial can have arbitrarily small relative degree:
interpolate parity on a sufficiently wide central $O(\sqrt m)$ window
of an $m$-cube and rescale to margin one, then hold $m$ fixed while
increasing the substituted block sizes.
Accordingly, this note supplies no symmetrization theorem and no
sample-complexity theorem for a general posterior.

\section{A Fourier cutoff from the degree bound}
The key estimate uses the concentration of a finite-dimensional
Schr\"odinger spectral subspace. We give the needed estimate directly.

\begin{lemma}[Weighted Fourier concentration]\label{lem:fourier}
Fix $0<\alpha<1$ and $\theta\in(0,\pi)$ with
$\alpha<1-\cos(\theta/2)$. There are $\kappa>0$ and $k_0$ such that,
for every $k\ge k_0$ and symmetric polynomial $P$ of degree
$r\le\alpha k$, the quantities
\[
 a_h=w_k(h)P(h),\qquad Q(t)=\sum_{h=0}^k a_h e^{ith},
 \qquad M=\|P\|_1
\]
satisfy
\begin{equation}\label{eq:fourier}
 \sup_{\theta\le|t|\le\pi}|Q(t)|\le M e^{-\kappa k}.
\end{equation}
No margin assumption is needed.
\end{lemma}
\begin{proof}
Let $\mathcal Q(z)=\sum_h w_k(h)P(h)z^h$, so that
$Q(t)=\mathcal Q(e^{it})$. The identity
\[
 \sum_{h=0}^k\binom kh(h)_jz^h=(k)_jz^j(1+z)^{k-j}
\]
for falling factorials shows that
$\mathcal Q(z)=(1+z)^{k-r}R(z)$, where $R$ is real and has degree
at most $r$. Put $m=k-r$ and $x=t/2$. Then
\[
 e^{-ikx}\mathcal Q(e^{2ix})
 =2^m\cos^m x\,e^{-irx}R(e^{2ix}).
\]
The last factor has frequencies $r,r-2,\ldots,-r$. Elementary
trigonometric identities imply that the real and imaginary parts of
the centered generating function, in one order according to the parity
of $r$, belong to the spaces
\[
 V_{m,r}=\{\cos^m x\,p(\sin x):\deg p\le r\},\qquad
 V_{m+1,r-1}.
\]
The second component is zero when $r=0$. The constant $2^m$ is
absorbed into the coefficients. Each real component has ordinary
$L^2(-\pi/2,\pi/2)$ norm at most $\sqrt\pi M$, because
$|Q(t)|\le M$.

For $\lambda\ge2$, let $p_n$ be orthonormal polynomials for the weight
$(1-y^2)^{\lambda-1/2}\,dy$ on $(-1,1)$. Integration by parts shows
that
\[
 L_\lambda p=-(1-y^2)p''+(2\lambda+1)yp'
\]
is symmetric for this weight. It preserves degree and has leading
coefficient $n(n+2\lambda)$ on a polynomial of degree $n$.
Orthogonality therefore gives $L_\lambda p_n=n(n+2\lambda)p_n$.
The functions $\phi_n(x)=\cos^\lambda x\,p_n(\sin x)$ are
orthonormal in $dx$, span $V_{\lambda,s}$ for $0\le n\le s$, and
satisfy, by direct differentiation,
\begin{equation}\label{eq:schrodinger}
 -\phi_n''+\lambda(\lambda-1)\sec^2x\,\phi_n
 =(n+\lambda)^2\phi_n.
\end{equation}
In both spaces above, $\lambda\ge m$ and $n+\lambda\le k$.

Choose
$\arccos(1-\alpha)<x_1<x_0=\theta/2$. For some $\delta>0$ and all
large $k$, the potential $V=\lambda(\lambda-1)\sec^2x$ and energy
$E=(n+\lambda)^2$ obey
\[
 V-E\ge m(m-1)\sec^2x-k^2\ge\delta k^2
 \qquad(|x|\ge x_1).
\]
Set $F(x)=k\sqrt{\delta/2}(|x|-x_1)_+$. Multiplying
\eqref{eq:schrodinger} by $e^{2F}\phi_n$ and integrating gives
\begin{equation}\label{eq:agmon}
 \int\left(|(e^F\phi_n)'|^2+
       (V-E-|F'|^2)e^{2F}|\phi_n|^2\right)dx=0.
\end{equation}
Boundary terms vanish since $\lambda\ge2$ for large $k$.
The Lipschitz weight is justified directly in the weak identity, or
by smoothing. Outside $[-x_1,x_1]$ the second coefficient is at least
$\delta k^2/2$; inside, $F=0$ and $V-E\ge-k^2$. Consequently
\[
 \|(e^F\phi_n)'\|_2\le k,\qquad
 \|e^F\phi_n\|_2^2\le1+2/\delta.
\]
The elementary one-dimensional bound
$|u(x)|^2\le2\|u\|_2\|u'\|_2$ for a function vanishing at the
endpoints now implies
\[
 |\phi_n(x)|\le C\sqrt{k}\,e^{-ck}
 \qquad(|x|\ge x_0),
\]
where $c=\sqrt{\delta/2}(x_0-x_1)>0$. Expand either real component
in its orthonormal basis and use Cauchy--Schwarz. There are at most
$k+1$ terms, so each component is bounded by
$Ck e^{-ck}$ times its $L^2$ norm in this region. The norm bound
$\sqrt\pi M$ proves \eqref{eq:fourier}, after decreasing $c$ to
absorb the polynomial factor and increasing $k_0$.
\end{proof}

The low-pass interpolation
\begin{equation}\label{eq:F}
 F(z)=\frac1{2\pi}\int_{-\theta}^{\theta}Q(t)e^{-itz}\,dt
\end{equation}
is entire and real on the real axis. Since $|Q(t)|\le M$, it satisfies
\begin{equation}\label{eq:Fbounds}
 |F(z)|\le M e^{\theta|\operatorname{Im}z|},\qquad
 |F(h)-a_h|\le M e^{-\kappa k}\quad(h\in\mathbb Z).
\end{equation}
The second inequality follows by subtracting the omitted part of the
Fourier inversion formula. A hard frequency cutoff suffices.

\section{Parity advantage forces sign transitions}
\begin{lemma}\label{lem:transitions}
For any $g_0,\ldots,g_k\in\{-1,1\}$,
\begin{equation}\label{eq:transition}
 \left|\sum_{h=0}^k(-1)^h w_k(h)g_h\right|
 \le\sum_{h:g_h\ne g_{h+1}}w_{k-1}(h).
\end{equation}
\end{lemma}
\begin{proof}
The alternating cumulative-binomial identity is
$\sum_{i=0}^h(-1)^i\binom{k}{i}=(-1)^h\binom{k-1}{h}$.
Summation by parts therefore expresses the sum on the left as
\[
 -2^{-k}\sum_{h=0}^{k-1}(g_{h+1}-g_h)(-1)^h\binom{k-1}{h}.
\]
Taking absolute values proves the result.
\end{proof}

We need a simple discrete logarithmic kernel. For $R>1$ define
\[
 L_R(j)=\left[\log\frac{R}{|j|+1}\right]_+,
 \qquad S_R=\sum_{j\in\mathbb Z}L_R(j).
\]
Stirling's formula, or comparison with an integral, gives
\begin{equation}\label{eq:SR}
 S_R=2R+O(\log R).
\end{equation}
If $F(j)$ and $F(j+1)$ have opposite signs, a real zero lies between
them. In Jensen's formula centered at an integer $h$, that zero
contributes at least $L_R(j-h)$ at radius $R$.

\section{Averaging Jensen's formula}
\begin{proof}[Proof of Theorem~\ref{thm:main}]
By \eqref{eq:criterion}, choose
$2\arccos(1-\alpha)<\theta<\pi\rho$.
Use $F$ and $\kappa$ from \eqref{eq:F}--\eqref{eq:Fbounds}.
If $\log M>\kappa k/2$, the conclusion is already true for large $k$.
We may therefore assume $M\le e^{\kappa k/2}$.

Fix a large constant $B$ and then a small constant $\varepsilon>0$,
and put $R=\varepsilon\sqrt{k}$. Throughout the central window
$|h-k/2|\le(B+2\varepsilon)\sqrt{k}$, binomial weights are bounded
below by $c_{B,\varepsilon}/\sqrt{k}$. Thus \eqref{eq:Fbounds} and the
margin assumption imply, for all large $k$,
\begin{equation}\label{eq:anchors}
 \sign F(h)=\sign P(h),\qquad |F(h)|\ge\tfrac12 w_k(h).
\end{equation}
Let $J$ be the set of sign transitions $j$ with
$|j-k/2|\le B\sqrt{k}$. Lemma~\ref{lem:transitions} and a binomial tail
bound imply
\begin{equation}\label{eq:Jmass}
 \sum_{j\in J}w_k(j)\ge\rho-\tau_B-o(1),
 \qquad \tau_B\longrightarrow0\quad(B\longrightarrow\infty).
\end{equation}
Here $w_{k-1}(j)/w_k(j)=2(k-j)/k=1+o(1)$ uniformly in the fixed central
window. For example one can take
$\tau_B=2e^{-2(B-o(1))^2}$ before absorbing the vanishing term.

For central $j$ and $|h-j|\le R$, adjacent binomial ratios give
\begin{equation}\label{eq:ratios}
 w_k(h)\ge\eta_{B,\varepsilon,k}w_k(j),\qquad
 \liminf_{k\to\infty}\eta_{B,\varepsilon,k}
 \ge e^{-C\varepsilon(B+\varepsilon)}.
\end{equation}
Indeed each adjacent log ratio is $O((B+\varepsilon)/\sqrt{k})$ in
this window, and there are at most $\varepsilon\sqrt{k}$ such steps.

Apply Jensen's formula to $F$ centered at each integer $h$ in the
anchor window $I=\{|h-k/2|\le(B+2\varepsilon)\sqrt{k}\}$. Keeping only
the zeros supplied by $J$, and using \eqref{eq:Fbounds}, gives
\[
 \sum_{j\in J}L_R(j-h)
 \le \log M-\log|F(h)|+\frac{2\theta R}{\pi}.
\]
Zeros on the integration circle are handled by a limiting radius.
Multiply by $w_k(h)$ and sum over $h\in I$. The right side is at most
\[
 \log M+O_{B,\varepsilon}(\log k)+\frac{2\theta R}{\pi}
\]
by \eqref{eq:anchors}. On the left, \eqref{eq:ratios} and
\eqref{eq:SR} give at least
\[
 \eta_{B,\varepsilon,k}S_R\sum_{j\in J}w_k(j).
\]
Every nonzero kernel term has its anchor in $I$, so no term is missing.
Combining the estimates yields
\begin{equation}\label{eq:Jensenfinal}
 \log M\ge
 2R\left[\eta_{B,\varepsilon,k}(\rho-\tau_B-o(1))
                 -\frac\theta\pi\right]-O_{B,\varepsilon}(\log k).
\end{equation}
Choose $B$ so that $\tau_B$ is much smaller than $\rho-\theta/\pi$,
then choose $\varepsilon$ so that the loss in \eqref{eq:ratios} is also
smaller than this gap. The bracket is then bounded below by a positive
constant. Since $R=\varepsilon\sqrt{k}$, the theorem follows.
\end{proof}

\begin{proposition}[Vanishing margin defects]\label{prop:stable}
Under the strict parameter condition of Theorem~\ref{thm:main}, its
conclusion remains valid along sequences $P_k$ for which, for some
fixed $a>0$,
\[
 \Prb(|P_k|<a)=o(1),\qquad c(P_k)\ge\rho-o(1).
\]
More precisely, $\liminf k^{-1/2}\log\|P_k\|_1>0$.
\end{proposition}
\begin{proof}
Use the same central windows and remove transitions incident to a
vertex with $|P_k|<a$. Their total central binomial weight is $o(1)$,
so \eqref{eq:Jmass} changes by $o(1)$. At the remaining endpoints the
low-pass interpolation preserves signs. Also omit bad anchor vertices
from the averaged Jensen inequality. The lost left-hand contribution
is at most
\[
 \sum_{h\text{ bad}}w_k(h)\sum_{j\in\mathbb Z}L_R(j-h)
 =S_R\Prb(h\text{ bad})=o(R).
\]
At every retained anchor, $|F(h)|\ge a w_k(h)/2$. Thus
\eqref{eq:Jensenfinal} changes only by $o(R)+O_a(1)$, proving the claim.
\end{proof}
If a symmetric $P_k$ is $o(1)$ in $L^2$ from a globally margin-one
function $Z_k$ of parity error at most $(1-\rho)/2$, then the proposition
applies with $a=1/2$. The approximating function $Z_k$ itself need not
be symmetric.

\begin{corollary}[Symmetry in one block]\label{cor:block}
Let $P$ be a globally margin-one polynomial on $k$ target coordinates,
invariant under permutations of a block $B$ of $n$ coordinates. Suppose
its degree in the coordinates of $B$ is at most $\alpha n$ and its
global parity advantage is at least $\rho$, with the strict parameter
condition of Theorem~\ref{thm:main}. Arbitrary dependence on the other
coordinates is allowed. For all sufficiently large $n$,
\[
 \|P\|_1\ge c\exp(c'\sqrt n),
\]
where the positive constants depend only on $\alpha,\rho$.
\end{corollary}
\begin{proof}
Choose $0<\rho_0<\rho$ close enough that the strict parameter condition
still holds. Condition on all coordinates outside $B$, and absorb
their parity into the sign of the resulting symmetric polynomial.
The conditional parity advantages lie in $[-1,1]$ and have mean at
least $\rho$. Thus a fraction at least
$(\rho-\rho_0)/(1-\rho_0)$ of fibers have advantage at least $\rho_0$.
Each such fiber has the same degree and margin bounds, so
Theorem~\ref{thm:main} applies uniformly. Average its conclusion over
these fibers.
\end{proof}

\section{A consequence for all-radius margin repairs}
We now define an auxiliary geometric problem independently of any
neural-network posterior. Fix $0<\beta\le1/2$ and let the active points
be $u\in\{-1,1\}^k$, with an integer radius $h\ge0$. For any finitely
supported real charge $q=(q_{u,h})$ of total mass zero, put
\begin{equation}\label{eq:cusp}
 \mathcal V_\beta(q)=
 -\sum_{u,h,v,j}q_{u,h}q_{v,j}\bigl(d_H(u,v)+h+j\bigr)^\beta.
\end{equation}
This quadratic form is nonnegative: use the integral representation of
$t^\beta$ and positivity of the kernels
$e^{-t d_H(u,v)}e^{-th}e^{-tj}$. Define the energy of a repair $g$ by
\[
 \mathcal E_\beta(g)=
 \sup_{q:\,\sum q=0,\ \mathcal V_\beta(q)>0}
 \frac{\left|\sum_{u,h}q_{u,h}g(u,h)\right|^2}
      {\mathcal V_\beta(q)},
\]
with infinite energy if a zero-variance charge has nonzero pairing.
This definition specifies the only norm property used below.

A head $P$ has $|P(u)|\ge1$. At each wrong point $u$, a repair $g_u$
is admissible if
\[
 |P(v)+g_u(v,h)|\ge1\quad\text{for all }v,h,
 \qquad \sign(P(u)+g_u(u,0))=\chi_k(u),
\]
and $g_u(v,h)\to0$ as $h\to\infty$ for every $v$. Intermediate signs
are unrestricted. The last condition supplies eventual agreement with
the original head and is an explicit hypothesis of this model.

\begin{lemma}[Contact drop with arbitrary intermediate sign changes]
\label{lem:contact}
Suppose $P(u)$ and $P(v)$ have the same sign, and
$D=2+|P(u)|-|P(v)|>0$. Every admissible repair at $u$ has energy $E$
satisfying
\[
 E\ge c_\beta\min\left\{
 \frac{D^2}{(d_H(u,v)+1)^\beta},\ D^{2-\beta}\right\}.
\]
\end{lemma}
\begin{proof}
The charge comparing $(v,h)$ and $(v,h+1)$ has variance
\[
 v_h=2(2h+1)^\beta-(2h)^\beta-(2h+2)^\beta
 \le C_\beta(h+1)^{\beta-2}.
\]
A sign change between these radii requires a jump of at least $2$, and
therefore energy at least $4/v_h$. Choose an integer
$R\le C_\beta(1+E^{1/(2-\beta)})$ large enough that this lower bound
exceeds $E$ for every $h\ge R$. Eventual agreement implies agreement
already at radius $R$. Comparing $(u,0)$ with $(v,R)$ now gives a repair
difference of absolute value at least $D$. Its variance is
\[
 2(d_H(u,v)+R)^\beta-(2R)^\beta
 \le C_\beta\bigl[(d_H(u,v)+1)^\beta+E^{\beta/(2-\beta)}\bigr].
\]
Thus $D^2$ is bounded by $C_\beta E$ times the last bracket. At least
one of its two terms accounts for half of $D^2$, proving the result.
\end{proof}

\begin{theorem}[Symmetric heads have diverging population repair cost]
\label{thm:repair}
Fix $\beta\in(0,1/2]$ and a degree fraction $\alpha$ satisfying
\eqref{eq:criterion} with $\rho=1/3$. For all sufficiently large $k$,
every symmetric head $P$ with degree at most $\alpha k$, global margin
one, and parity error at most $1/3$, and every family of admissible
repairs, satisfy
\[
 \E_u\left[\mathbf1_{\{\sign P(u)\ne\chi_k(u)\}}
                  \mathcal E_\beta(g_u)\right]
 \ge \frac{c_\beta}{(k+1)^\beta}\exp(c\sqrt{k}),
\]
where $c>0$ may depend on $\alpha,\beta$.
\end{theorem}
\begin{proof}
Both signs occur, since a constant sign has parity error $1/2$.
Let $b_+$ and $b_-$ be the smallest positive and negative absolute head
values, and define
\[
 s=(b_++b_-)/2\ge1,\qquad t=(b_+-b_-)/2,\qquad Q=(P-t)/s.
\]
Then $Q$ is symmetric, has the same signs and degree bound, satisfies
$|Q|\ge1$, and has a contact at each of $+1$ and $-1$.
For every wrong $u$, choose a same-sign contact $v$. The contact drop
for the original head is
\[
 D=2+|P(u)|-b_{\sign P(u)}=2+s(|Q(u)|-1)\ge1+|Q(u)|.
\]
Lemma~\ref{lem:contact}, using $d_H(u,v)\le k$, gives, with
$p=2-\beta$,
\[
 \mathcal E_\beta(g_u)\ge
 \frac{c_\beta}{(k+1)^\beta}|Q(u)|^p.
\]
Since $\deg Q<k$, parity orthogonality gives $\E[\chi_k Q]=0$.
Consequently the wrong and correct sides have equal absolute first
moments, and
\[
 \E[\mathbf1_{\rm wrong}|Q|]=\tfrac12\|Q\|_1.
\]
H\"older's inequality implies
$\E[\mathbf1_{\rm wrong}|Q|^p]\ge2^{-p}\|Q\|_1^p$.
Apply Theorem~\ref{thm:main} to $Q$ and average the preceding energy
bound.
\end{proof}

The same proof applies to a head symmetric in a block of $n$ target
coordinates, with block degree at most $\alpha n$: replace the final
application of Theorem~\ref{thm:main} by Corollary~\ref{cor:block}.
Parity orthogonality still holds because the degree in that block is
less than $n$, regardless of the degree in the other coordinates.
The resulting lower bound is
$c(k+1)^{-\beta}\exp(c'\sqrt n)$. In particular it diverges
exponentially in $\sqrt k$ when the block contains a fixed positive
fraction of the target coordinates.

\begin{corollary}[Only the sign pattern needs block symmetry]
\label{cor:signblock}
Let $P$ have global margin one on $k$ target coordinates, and suppose
$\sign P$ is invariant under permutations of a block of $m$ coordinates.
The amplitudes and the dependence on coordinates outside that block may
be arbitrary. If the block degree is at most $\alpha m$, the parity
error is at most $0.34$, and
\[
 0<\alpha<1-\cos(4\pi/25)=0.123693319956\ldots,
\]
then every family of admissible repairs satisfies, for sufficiently
large $m$,
\[
 \E_u[\mathbf1_{\rm wrong}\mathcal E_\beta(g_u)]
 \ge c(m+1)^{-\beta}\exp(c'\sqrt m).
\]
The constants may depend on the strict gap in this degree condition.
\end{corollary}
\begin{proof}
Choose $0<\rho_0<0.32$ sufficiently close to $0.32$ that
$\alpha<1-\cos(\pi\rho_0/2)$. Condition on coordinates outside
the block. Since the mean conditional parity advantage is at least
$0.32$, a fraction at least $(0.32-\rho_0)/(1-\rho_0)>0$ of the
fibers have advantage at least $\rho_0$. Fix such a fiber and absorb the fixed outside
parity into its target $f$. Write $s=\sign P$, $E=\{fs=-1\}$,
$C=\{fs=1\}$, $R=\Prb(E)\le(1-\rho_0)/2$, and $p=2-\beta$. Define
\[
 \Phi_s(P)=\E_{u}\mathbf1_E(u)
   \max_{v\in C:\,s(v)=s(u)}
   \frac{[s(u)(P(u)-P(v))]_+^p}{(1+d_H(u,v))^\beta}.
\]
This functional is convex in $P$ and invariant under permutations of
the block. Averaging $P$ over those permutations gives a symmetric
polynomial $\overline P$ with the same signs, margin, and degree bound,
and $\Phi_s(\overline P)\le\Phi_s(P)$. The contact-drop lemma implies
that the original repair cost on this fiber is at least
$c_\beta\Phi_s(P)$.

For completeness, $\Phi_s$ controls the first moment uniformly over
this sign cone. Put $c_0=1-2R$ and
$b_\sigma=\min_{v\in C:s(v)=\sigma}|\overline P(v)|$.
Each phase has correct mass minus wrong mass $c_0/2$. With
$D_u=(|\overline P(u)|-b_{s(u)})_+$ on $E$, parity orthogonality gives
\[
 \frac{c_0}{2}(b_++b_-)\le\E[\mathbf1_E D_u],\qquad
 \|\overline P\|_1\le\frac2{c_0}\E[\mathbf1_E D_u].
\]
Consequently H\"older's inequality yields
\[
 \|\overline P\|_1
 \le\frac2{c_0}R^{1-1/p}(m+1)^{\beta/p}
       \Phi_s(\overline P)^{1/p}.
\]
Theorem~\ref{thm:main}, with advantage $\rho_0$, now gives
$\Phi_s(P)\ge c(m+1)^{-\beta}\exp(c'\sqrt m)$.
Averaging over the favorable fibers proves the claim. Restricting
contact tests to one fiber is valid for every original repair, so no
symmetry of the repairs or of the original amplitudes was imposed.
\end{proof}

This result allows every intermediate sign change and every choice of
repair values. Its eventual-agreement condition and the stated symmetry
of the head or its sign pattern remain assumptions of the statement. In particular,
neither the rotational symmetry of a prior nor the radial form of a
kernel makes a head fitted to an individual training set symmetric.
A general posterior lower bound requires an additional argument beyond
this note.

\begin{thebibliography}{9}
\bibitem{odonnell} R.~O'Donnell, \emph{Analysis of Boolean Functions},
Cambridge University Press, 2014. Boolean Fourier expansions.
\bibitem{ahlfors} L.~V.~Ahlfors, \emph{Complex Analysis}, 3rd ed.,
McGraw--Hill, 1979. Jensen's formula.
\end{thebibliography}
\end{document}
